% D29 PR-AUC re-reading, generated by scripts/analysis/d29_read.py (do not hand-edit the numbers;
% re-run the reader instead). Drop-in replacement for Table tab:d29pr and Table tab:d29prc of
% sections/app-c-robust.tex, plus the appendix paragraph and the main-text sentence below.
% Columns are sized for the TMLR text width at \scriptsize; booktabs, as the rest of the appendix.

% ===== Table 1 of 2: C1 =====
\begin{table}[!htbp]
\caption{D29 PR-AUC readings of C1 (descriptive; existing records, test split at the ROC-selected
epoch): \lastarm$-$\lparm{} in average precision (AP) and in the trapezoid PR-AUC of the release
($\times100$). No practical-importance threshold is applied to a PR metric: the $\pm3$-point margin of
our rules is a ROC margin, no PR-AUC margin was pre-registered, and the PR baselines here span
0.89 to 42.21 points. Each row gives the two arm means, the difference with its
95\% $t$-interval over the units (5
seeds on DL's split, 5 rotation means of 3 seeds elsewhere) and the change relative to the \lparm{}
baseline, which should be read together with that baseline: a percentage on a near-prevalence
baseline is
unstable. A dagger ($\dagger$) marks the
10 intervals that exclude 0, the only inference
claim made here. The ROC column is the unchanged ROC reading with its consistency
label (F, S, N: Table~\ref{tab:vocab}). Heads: raw = linear head, MLP = shipped MLP head (fixed
recipe),
std = \texttt{linear\_z}, tuned raw / tuned std = the D28 probes; run tags as in
Table~\ref{tab:glossary}.}
\label{tab:d29pr}
\centering\scriptsize
\setlength{\tabcolsep}{2.5pt}
\begin{tabular}{llllllll}
\toprule
model & head (tags) & ROC & AP \lparm{} / \lastarm{} & AP difference [95\% $t$] & AP rel. & trapezoid & rel.\\
\midrule
\multicolumn{8}{l}{\emph{DL's split}}\\
BIOT & raw (P1E) & $-$0.60 N & 1.94 / 1.91 & $-$0.03 [$-$0.41, +0.36] & $-$1.4\% & $-$0.03 & $-$1.5\%\\
CBraMod & raw (P1E) & +0.01 N & 42.21 / 42.41 & +0.20$^{\dagger}$ [+0.08, +0.32] & +0.5\% & +0.20$^{\dagger}$ & +0.5\%\\
BIOT & MLP (P1E) & $-$1.19 N & 3.14 / 2.67 & $-$0.47 [$-$1.22, +0.27] & $-$15.0\% & $-$0.47 & $-$15.2\%\\
CBraMod & MLP (P1E) & +0.03 N & 39.21 / 38.58 & $-$0.63 [$-$2.85, +1.60] & $-$1.6\% & $-$0.63 & $-$1.6\%\\
BIOT & std (P1Z) & +0.18 N & 20.24 / 21.92 & +1.68 [$-$3.40, +6.76] & +8.3\% & +1.68 & +8.4\%\\
CBraMod & std (P1Z) & +0.18 N & 34.12 / 34.70 & +0.58$^{\dagger}$ [+0.30, +0.86] & +1.7\% & +0.60$^{\dagger}$ & +1.8\%\\
BIOT & tuned raw (P1G) & $-$0.08 N & 12.63 / 13.02 & +0.39 [$-$0.31, +1.09] & +3.1\% & +0.39 & +3.1\%\\
CBraMod & tuned raw (P1G) & $-$0.54 N & 37.29 / 35.92 & $-$1.37 [$-$4.66, +1.92] & $-$3.7\% & $-$1.38 & $-$3.7\%\\
BIOT & tuned std (P1G) & $-$0.09 N & 20.18 / 18.44 & $-$1.74 [$-$9.09, +5.62] & $-$8.6\% & $-$1.75 & $-$8.8\%\\
CBraMod & tuned std (P1G) & +0.12 N & 35.28 / 35.85 & +0.57$^{\dagger}$ [+0.10, +1.04] & +1.6\% & +0.57$^{\dagger}$ & +1.6\%\\
\multicolumn{8}{l}{\emph{CHB-MIT}}\\
BIOT & raw (S1E+S1M) & +1.91 N & 0.92 / 1.56 & +0.64 [$-$1.00, +2.28] & +70.0\% & +0.65 & +72.3\%\\
CBraMod & raw (S1E+S1M) & $-$0.02 N & 23.02 / 23.29 & +0.27 [$-$0.33, +0.87] & +1.2\% & +0.27 & +1.2\%\\
BIOT & std (S1Z) & +0.11 N & 6.80 / 7.19 & +0.39 [$-$0.39, +1.17] & +5.7\% & +0.38 & +5.7\%\\
CBraMod & std (S1Z) & +0.41 N & 28.75 / 29.58 & +0.82 [$-$0.19, +1.84] & +2.9\% & +0.82 & +2.9\%\\
BIOT & tuned raw (S1G) & $-$0.51 N & 11.51 / 11.61 & +0.09 [$-$0.92, +1.11] & +0.8\% & +0.11 & +1.0\%\\
CBraMod & tuned raw (S1G) & +0.66 N & 19.39 / 19.75 & +0.36 [$-$1.58, +2.29] & +1.8\% & +0.37 & +1.9\%\\
BIOT & tuned std (S1G) & $-$0.39 N & 8.67 / 8.57 & $-$0.10 [$-$1.55, +1.35] & $-$1.2\% & $-$0.10 & $-$1.1\%\\
CBraMod & tuned std (S1G) & $-$4.42 S & 25.60 / 25.36 & $-$0.24 [$-$3.63, +3.14] & $-$0.9\% & $-$0.23 & $-$0.9\%\\
\multicolumn{8}{l}{\emph{Siena}}\\
BIOT & raw (S3E+S3M) & +9.23 F & 1.65 / 2.20 & +0.55$^{\dagger}$ [+0.14, +0.96] & +33.2\% & +0.54$^{\dagger}$ & +36.2\%\\
CBraMod & raw (S3E+S3M) & +0.27 N & 10.50 / 10.84 & +0.34 [$-$0.83, +1.50] & +3.2\% & +0.30 & +3.0\%\\
BIOT & std (S3Z) & +1.65 N & 6.37 / 6.55 & +0.18$^{\dagger}$ [+0.01, +0.35] & +2.8\% & +0.19$^{\dagger}$ & +3.0\%\\
CBraMod & std (S3Z) & $-$0.26 N & 9.05 / 9.09 & +0.04 [$-$0.98, +1.06] & +0.4\% & +0.00 & +0.0\%\\
BIOT & tuned raw (S3G) & +0.18 N & 5.71 / 5.25 & $-$0.46 [$-$1.32, +0.40] & $-$8.0\% & $-$0.46 & $-$8.6\%\\
CBraMod & tuned raw (S3G) & +0.40 N & 13.33 / 13.14 & $-$0.19 [$-$2.55, +2.17] & $-$1.4\% & $-$0.17 & $-$1.3\%\\
BIOT & tuned std (S3G) & $-$0.17 N & 6.99 / 7.42 & +0.43 [$-$0.70, +1.56] & +6.1\% & +0.43 & +6.5\%\\
CBraMod & tuned std (S3G) & +0.57 N & 9.00 / 9.33 & +0.33 [$-$1.00, +1.65] & +3.7\% & +0.29 & +3.3\%\\
\multicolumn{8}{l}{\emph{Siena filtered}}\\
BIOT & raw (S3FE+S3FM) & +4.11 F & 6.23 / 6.05 & $-$0.18 [$-$4.03, +3.67] & $-$2.9\% & $-$0.18 & $-$3.1\%\\
CBraMod & raw (S3FE+S3FM) & +0.14 N & 27.28 / 28.23 & +0.94 [$-$1.99, +3.88] & +3.5\% & +0.97 & +3.6\%\\
BIOT & std (S3FZ) & +1.93 N & 5.49 / 5.77 & +0.28 [$-$0.16, +0.71] & +5.1\% & +0.28 & +5.2\%\\
CBraMod & std (S3FZ) & +0.34 N & 10.91 / 11.09 & +0.18 [$-$0.57, +0.93] & +1.6\% & +0.20 & +1.9\%\\
BIOT & tuned raw (S3FG) & +0.69 N & 8.84 / 11.08 & +2.23 [$-$4.68, +9.15] & +25.3\% & +2.23 & +26.7\%\\
CBraMod & tuned raw (S3FG) & +0.18 N & 27.08 / 28.57 & +1.49 [$-$0.59, +3.56] & +5.5\% & +1.54 & +5.8\%\\
BIOT & tuned std (S3FG) & +0.71 N & 7.56 / 7.90 & +0.34 [$-$0.84, +1.52] & +4.5\% & +0.26 & +3.6\%\\
CBraMod & tuned std (S3FG) & $-$1.30 N & 11.63 / 11.19 & $-$0.44 [$-$1.61, +0.74] & $-$3.8\% & $-$0.41 & $-$3.6\%\\
\bottomrule
\end{tabular}
\end{table}
% src: artifacts/analysis/D29/pr_auc.md section 1 (AP and trapezoid rows) and 1b (ROC reference);
% artifacts/analysis/D29/d29.json "pr"."C1"; the reading itself: d29.json "pr"."reading"

% ===== Table 2 of 2: C2/C3 and the noise arm =====
\begin{table}[!htbp]
\caption{D29 PR-AUC readings of C2/C3 and of the noise arm (descriptive), in average precision
($\times100$), with no threshold applied: spans over the curve points ($N$ 896--4174 and $P$ 2/5/10 at
$N$ 2048; 4 points for the CBraMod control). \emph{curve $-$ LP} and \emph{vs \lastnoise} give the
range of the per-point differences and, in the next column, of the changes relative to the baseline arm;
\emph{excl.} counts the points whose 95\% $t$-interval excludes 0 (the only inference claim), and the
number in brackets the points whose wider Nadeau-Bengio interval also does (that interval is defined for
the rotation settings only). Units: 5 rotation means of 3
seeds on Siena, 5 seeds on DL's split for BIOT; CBraMod on DL's split has a single curve seed, so no
interval exists for that row. DL's split: as in Table~\ref{tab:c2c3}; Siena: BIOT the D29 runs, CBraMod
the fixed-compute control at seed $S_r$ (linear head only, no noise arm). The trapezoid PR-AUC gives the
same picture throughout (it differs from AP by at most 0.12 points in any contrast); per-point
numbers and intervals for both metrics are in \texttt{artifacts/analysis/D29/pr\_auc.md}.}
\label{tab:d29prc}
\centering\scriptsize
\setlength{\tabcolsep}{2pt}
\begin{tabular}{llllllll}
\toprule
setting & model, head & LP & curve & curve $-$ LP & rel. & vs \lastnoise & rel.\\
\midrule
DL's split & BIOT, raw & 5.56 & 5.07--5.60 & $-$0.49 to +0.04 (0/7) & $-$8.8 to +0.7\% & -- & --\\
DL's split & CBraMod, raw & 41.68 & 43.96--44.12 & +2.28 to +2.44 (--) & +5.5 to +5.9\% & -- & --\\
Siena & BIOT, raw & 1.64 & 2.59--2.70 & +0.95 to +1.05 (6/7 [0]) & +57.7 to +64.1\% & +0.35 to +0.50 (7/7 [7]) & +15.7 to +22.9\%\\
Siena & BIOT, std & 6.24 & 6.68--6.78 & +0.44 to +0.54 (0/7 [0]) & +7.1 to +8.6\% & +0.17 to +0.43 (0/7 [0]) & +2.5 to +6.7\%\\
Siena & CBraMod, raw & 12.85 & 14.86--14.89 & +2.01 to +2.04 (0/4) & +15.6 to +15.9\% & -- & --\\
Siena filt. & BIOT, raw & 5.77 & 5.53--5.74 & $-$0.24 to $-$0.03 (0/7 [0]) & $-$4.1 to $-$0.6\% & +0.58 to +0.77 (0/7 [0]) & +11.6 to +15.6\%\\
Siena filt. & BIOT, std & 5.66 & 5.64--5.73 & $-$0.02 to +0.08 (0/7 [0]) & $-$0.3 to +1.4\% & $-$0.11 to +0.06 (0/7 [0]) & $-$1.9 to +1.1\%\\
Siena filt. & CBraMod, raw & 28.01 & 32.93--33.18 & +4.92 to +5.18 (0/4) & +17.6 to +18.5\% & -- & --\\
\bottomrule
\end{tabular}
\end{table}
% src: artifacts/analysis/D29/pr_auc.md sections 2 (DL's split), 3 (vs LP), 4 (noise arm) and 5
% (CBraMod control); artifacts/analysis/D29/d29.json "pr"."C2C3_dl_split" / "C2C3_rotations"

% ===== Appendix paragraph (replaces the "PR-AUC (descriptive)" paragraph of app-c-robust.tex) =====
\emph{PR-AUC (descriptive).} We read C1, C2 and C3 again in average precision (AP) and in the trapezoid
PR-AUC of the release, and we do not carry the $\pm3$-point margin of our ROC rules over to them: that
margin was fixed for ROC-AUC, no PR-AUC margin was pre-registered, and AP sits near the test prevalence
(1.49\% on DL's split, 0.16--0.66\% on the CHB-MIT rotations, 0.43--0.99\% on Siena), so one fixed
margin would mean very different things across these settings. We therefore report, for each contrast,
the two arm means, the difference with its 95\% $t$-interval over the units, the change relative to the
baseline arm, and whether the interval excludes 0; we attach no practical-importance label, and a
relative change should always be read next to the baseline printed beside it, since a percentage on a
near-prevalence baseline is unstable. In Table~\ref{tab:d29pr}, 10 of 68 C1 contrasts
have an interval that excludes 0, all of them positive; the largest relative change is Siena BIOT with
the raw head, where AP moves from 1.65 to 2.20 (+0.55
[+0.14, +0.96], about +33\% of that 1.65-point baseline) --- the same
comparison that reads a consistent difference of +9.23 ROC points. On C2/C3
(Table~\ref{tab:d29prc}) the Siena BIOT raw-head curve sits +0.95 to +1.05 AP
points above LP at every subset and +0.35 to +0.50 above its own \lastnoise{} arm; the EEG-versus-noise
difference is the one that also survives the wider Nadeau-Bengio interval at every subset, while the
gain over LP does not, and \lastnoise{} itself reaches 2.19 to
2.24 AP against LP 1.64, so most of the gain over LP
is reproduced by perturbing with noise rather than by the EEG data. In every other row of that table
the intervals include 0 on both PR metrics, or (CBraMod on DL's split, one curve seed) no interval
exists. None of this changes a ROC reading: the D29 noise rule
still reads ``no difference demonstrated'' on ROC-AUC, where the same contrast has means of about $+3$
to $+4$ points with intervals that include 0.
% src: artifacts/analysis/D29/README.md "PR-AUC readings" (the reading, prevalence, the C1 rows whose
% interval excludes 0, the rotation span table and its reading paragraph); pr_auc.md sections 1-6

% ===== Main-text sentence (one sentence; place where the PR-AUC aside is made) =====
We also read precision-recall performance without the $\pm3$-point margin of our ROC rules, which was
never defined for it: the adaptation differences that ROC-AUC resolves are small in absolute average
precision but large next to a near-prevalence baseline (Siena BIOT, raw head: 1.65 to
2.20 AP, +0.55 [+0.14, +0.96], about +33\% of the
\lparm{} baseline), and we report them as measured differences with intervals rather than labelling
them against a ROC-scale threshold (Appendix~\ref{app:d29}).
% src: artifacts/analysis/D29/README.md "PR-AUC readings"; the ROC readings are unchanged
